\section{RL for LLM Reasoning}

\subsection{Más allá de las preferencias: recompensas verificables}

En dominios como las matemáticas y la programación, la corrección de una respuesta puede \textbf{verificarse automáticamente}. En ese caso no se necesitan preferencias humanas: la corrección puede usarse directamente como recompensa para el entrenamiento con RL.

\begin{importantbox}{El auge de RL for Reasoning}
Modelos como DeepSeek-R1 mostraron que el entrenamiento con RL permite a un LLM aprender capacidades como la autocorrección durante el razonamiento y la cadena de pensamiento. Esta es otra dirección de aplicación importante de RL en los LLM, además de la alineación.
\end{importantbox}

\subsection{Diseño de recompensas verificables}

Las señales habituales en tareas verificables incluyen:
\begin{itemize}
    \item \textbf{Consistencia lógica}: determinar automáticamente si la solución de un problema matemático cumple la lógica de sus pasos.
    \item \textbf{Razonamiento sintético}: revisar la cadena de razonamiento con un symbolic executor; la reward se retroalimenta directamente al token.
    \item \textbf{Verificación generativa}: usar un verifier LM para calificar las respuestas candidate, como auxiliary reward.
\end{itemize}

\begin{table}[H]
\centering
\begin{tabular}{p{0.3\textwidth} p{0.35\textwidth} p{0.32\textwidth}}
\toprule
Categoría de tarea & Señal verificable & Métrica típica \\
\midrule
Razonamiento matemático & steps correctness & chain-of-thought accuracy, full solution correctness \\
Programación & unit test pass rate & pass@k, runtime exceptions \\
Preguntas y respuestas de conocimiento & factuality consistency & TruthfulQA, F1 against evidence \\
\bottomrule
\end{tabular}
\caption{Tareas y métricas que cubren las recompensas verificables}
\end{table}

\begin{knowledgebox}{La función del verifier autosupervisado}
Cuando la comparación humana no es viable, introducir un verifier autosupervisado (como un program executor o un verifier LM) permite generar la reward signal y asegura el ciclo cerrado del reasoning RL.
\end{knowledgebox}

\subsection{Experimentos y evaluación}

Trabajos como DeepSeek-R1 suelen usar como baseline: Instruction tuned LM + RLHF contrast, DPO y el RL for reasoning pipeline. Los métodos de verificación incluyen la generación con LM, los prompts de varios turnos y la revisión con un external verifier.

\begin{importantbox}{Lo que revela el Reasoning RL}
\textit{``RL proporciona al modelo un mecanismo de autocorrección, en lugar de solo reforzar la imitación''}: esta retroalimentación hace que el model, ante razonamientos de cadena larga o multi-hop question, esté más dispuesto a traverse deeper reasoning tree.
\end{importantbox}

En la evaluación también debe considerarse la inference latency, porque el reasoning RL suele requerir rollouts y un verifier; al desplegarlo hay que sopesar la experiencia en real-time.

\subsection{Resumen de la sección}

En los LLM, RL no solo se usa para la alineación (RLHF/DPO), sino también para el reasoning: mediante recompensas verificables y un verifier pipeline, el modelo obtiene supervisión para autocorregirse; en los experimentos hay que considerar a la vez el trade-off entre correctness y latency.
